% Simple LaTeX report for CA20
\documentclass{article}
\usepackage{graphicx}
\usepackage{hyperref}
\title{CA20: Lagrangian-constrained Policy Optimization — Short Report}
\author{Student Name}
\date{\today}
\begin{document}
\maketitle
\begin{abstract}
This short report documents a minimal experiment using a simple MLP policy trained with a Lagrangian relaxation on a synthetic dataset.
\end{abstract}
\section{Method}
We implement a Gaussian MLP policy with a state-independent log-std and a value network for baseline advantage estimation. The Lagrangian objective is:\\
$L = $ policy loss $ + \mu (C - c)$\\
\section{Experiments}
Run the debug config and save figures:\\
\texttt{python -c "from paperAssignments.Assignments1_50.CA20.src import train, config; cfg=config.Config(); cfg.epochs=2; cfg.batch_size=32; train.train(cfg)"}
\section{Results}

This section contains example placeholders for tables and figures. Replace the placeholders with actual results produced by running the experiments described in Section~\ref{sec:experiments}.

\subsection{Learning Curves}
\begin{figure}[h]
  \centering
  % Replace with your reward curve: e.g. notebooks/pictures/reward_curve.png
  \includegraphics[width=0.7\textwidth]{notebooks/pictures/reward_curve.png}
  \caption{Average reward over training steps for a representative seed (placeholder).}
  \label{fig:reward}
\end{figure}

\begin{figure}[h]
  \centering
  % Replace with your constraint violation curve: e.g. notebooks/pictures/constraint_curve.png
  \includegraphics[width=0.7\textwidth]{notebooks/pictures/constraint_curve.png}
  \caption{Mean constraint violation over training steps (placeholder).}
  \label{fig:constraint}
\end{figure}

\subsection{Summary Tables}
\begin{table}[h]
  \centering
  \begin{tabular}{lrrrr}
    \hline
    Method & Mean return & Std return & Mean constraint & Std constraint\\
    \hline
    Baseline & -- & -- & -- & -- \\
    Lagrangian ($c=0.05$) & -- & -- & -- & -- \\
    Lagrangian ($c=0.1$) & -- & -- & -- & -- \\
    \hline
  \end{tabular}
  \caption{Summary statistics across seeds (placeholders).}
  \label{tab:summary}
\end{table}

\section{Discussion}

Interpret the results and discuss potential failure modes. For example:

- How does the multiplier evolve over time and does it stabilise? Compare final multiplier values across runs.
- Is there a visible trade-off between reward and constraint satisfaction in the sweep over $c$?
- What hyperparameters (learning rates, lambda lr, clipping) affected stability the most?

\section{Reproducibility and Build Instructions}
\label{sec:repro}

To reproduce these results locally:

\begin{enumerate}
  \item Create and activate a Python 3.10+ virtualenv.
  \item Install dependencies: \texttt{pip install -r requirements.txt}.
  \item Run the debug experiment (short run):
  \begin{verbatim}
python -c "from paperAssignments.Assignments1_50.CA20.src import train, config; cfg=config.Config(); cfg.epochs=2; cfg.batch_size=32; train.train(cfg)"
  \end{verbatim}
  \item Run the notebook `notebooks/demo.ipynb` to generate example figures saved to `notebooks/pictures/`.
  \item Compile this report to PDF (requires a LaTeX toolchain):
  \begin{verbatim}
pdflatex report.tex
bibtex report || true
pdflatex report.tex
pdflatex report.tex
  \end{verbatim}
\end{enumerate}

\section{Conclusion}

We presented a compact, reproducible scaffold for studying Lagrangian-constrained optimization in simple RL settings. The codebase is intentionally minimal and designed for pedagogy and extension; future work includes integrating environment-based episodic evaluation, richer policy classes, and more extensive hyperparameter sweeps.

\section*{Acknowledgements}

This scaffold was developed for course CA20. Please cite this repository if you use it for course work or research.

\section*{Data and Code Availability}

All code is included in this repository. Example scripts and notebooks save figures and checkpoints under `notebooks/pictures/` and `outputs/` respectively. Tests are located in `tests/` and can be run with `pytest`.

\bibliographystyle{plain}
\bibliography{refs}
\end{document}
